\section{If-Then Statements and Proof Strategy}
\chapterline

In the previous chapter, we treated definitions as precise classification tools: given an object, they tell us whether it belongs to a class. In mathematics, we connect these conditions into chains of consequence:
\[
P \implies Q.
\]
We call this an \textbf{implication}---a conditional promise that whenever condition $P$ is satisfied, outcome $Q$ is guaranteed to follow. The statement $P$ is called the \textbf{hypothesis} (the starting assumption), and $Q$ is called the \textbf{conclusion} (the promised result).

When you are asked to prove an implication, beginners often make an understandable initial mistake: they look at $P$ and wonder whether it is actually true in the real world. But an implication never asserts that $P$ is true! It makes a conditional promise:
\begin{center}
\itshape ``In any mathematical scenario where $P$ happens to hold, $Q$ is guaranteed to hold as well.''
\end{center}
Think of an airline boarding pass: ``If you hold a ticket for Seat 14B, you may board the plane.'' The airline is not claiming that you hold such a ticket. They are stating a binding rule: hand them a ticket for 14B, and a seat on the plane is guaranteed.

\subsection{Implication as geometric containment}

The cleanest way to visualize an implication is through the geometry of sets. Let $U$ be the universe of all possible scenarios, let $P$ be the set of scenarios where the hypothesis holds, and let $Q$ be the set of scenarios where the conclusion holds.

The assertion $P \implies Q$ means precisely that \textbf{$P$ is a subset of $Q$} ($P \subseteq Q$).

\begin{center}
\begin{tikzpicture}[
    font=\small,
    universe/.style={draw=gray!50, fill=gray!2, rounded corners=4pt, thick},
    concl/.style={draw=sourceblue!80, fill=sourceblue!12, rounded corners=12pt, thick},
    hypo/.style={draw=targetgreen!80, fill=targetgreen!25, rounded corners=8pt, thick}
]

% Universe Box
\draw[universe] (0, 0) rectangle (12.5, 4.4);
\node[anchor=north west, font=\footnotesize\scshape, text=gray!70] at (0.2, 4.2) {Universe of All Scenarios $U$};

% Conclusion Region Q
\draw[concl] (1.0, 0.6) rectangle (8.0, 3.6);
\node[font=\footnotesize\bfseries, text=sourceblue!90!black] at (6.2, 3.2) {Conclusion $Q$ Holds};

% Hypothesis Region P (nested entirely inside Q)
\draw[hypo] (1.5, 1.0) rectangle (4.5, 2.8);
\node[font=\footnotesize\bfseries, text=targetgreen!90!black] at (3.0, 2.4) {Hypothesis $P$ Holds};
\node[font=\scriptsize, text=targetgreen!90!black, align=center] at (3.0, 1.6) {Region (1):\\$P \land Q$\\Direct Verification};

% Vacuous region inside Q but outside P
\node[font=\scriptsize, text=sourceblue!80!black, align=center] at (6.2, 1.7) {Region (2):\\$\neg P \land Q$\\\textbf{Vacuously True}\\Contract honored};

% Vacuous region outside both P and Q
\node[font=\scriptsize, text=gray!80!black, align=center] at (10.2, 2.8) {Region (3):\\$\neg P \land \neg Q$\\\textbf{Vacuously True}\\Contract unbroken};

% Forbidden breach region (highlighted outside)
\draw[draw=trapred!60, fill=trapred!10, rounded corners=3pt, thick, dashed] (8.5, 0.6) rectangle (12.0, 1.8);
\node[font=\scriptsize\bfseries, text=trapred!90!black, align=center] at (10.25, 1.2) {Forbidden Breach:\\$P \land \neg Q$\\\textbf{MUST BE EMPTY}};

\end{tikzpicture}
\end{center}

This geometry reveals how truth behaves:
\begin{itemize}
\item \textbf{Region 1 ($P$ is true, $Q$ is true)}: You stand inside $P$. Because the boundary of $P$ is nested entirely inside $Q$, you are automatically inside $Q$. For example, if $n$ is a multiple of $4$, then $n$ is automatically even. Both hypothesis and conclusion hold; the promise is fulfilled.
\item \textbf{Region 2 ($\neg P$ holds, $Q$ is true)}: You stand inside $Q$, but outside $P$. For example, the number $6$ is even, but it is not a multiple of $4$. Has the promise ``if a multiple of $4$, then even'' been broken? Not at all. The promise said nothing about what happens to numbers that are not multiples of $4$.
\item \textbf{Region 3 ($\neg P$ holds, $\neg Q$ holds)}: You stand outside both $P$ and $Q$. For example, $7$ is neither a multiple of $4$ nor even. Again, no promise was broken.
\item \textbf{The Forbidden Region ($P$ is true, $\neg Q$ holds)}: This would be a point inside $P$ but outside $Q$---for example, an integer that is a multiple of $4$ yet odd. If even a single such point existed, the implication $P \implies Q$ would be false. Saying that $P \implies Q$ is true means this region contains no points.
\end{itemize}

\subsection{Vacuous truth: promises about empty conditions}

The fact that $P \implies Q$ is considered true whenever $P$ is false is known as \textbf{vacuous truth}---a conditional promise that is declared true simply because its starting condition is impossible or never satisfied. Beginners often find this unsettling at first encounter. Consider the statement:
\begin{quote}
``If an integer is both a multiple of $4$ and an odd number, then that integer is purple.''
\end{quote}
Is this statement true? In mathematical logic, the answer is an unqualified \textbf{yes}. 

Why? Because an integer that is simultaneously a multiple of $4$ and odd does not exist. The hypothesis describes the empty set $\emptyset$---the unique set containing zero elements. To catch someone breaking an if-then promise, you must bring forward a real counterexample that fulfills premise $P$ while violating conclusion $Q$. Because no one can ever produce an odd multiple of $4$, no one can ever catch the promise in a breach.

\begin{recoverybox}
\textbf{Pause and think:} Suppose an instructor announces: ``If you score $105\%$ on the final exam, you will receive an A+.'' You take the exam, earn $88\%$, and receive a B. Did the instructor break their promise or lie to you?

\emph{Answer:} No. The promise was conditioned on scoring $105\%$. Because you never fulfilled that condition, the instructor had no obligation to award the A+, and no breach of contract occurred. In mathematics, unbroken promises are classified as true.
\end{recoverybox}

\begin{historybox}
Around 300 BCE, two Greek philosophers of the Megarian school---an ancient philosophical circle centered in Megara devoted to logic puzzles and paradoxes---sparked a debate over conditionals that lasted for more than two millennia. Philo of Megara and his teacher Diodorus Cronus asked a direct question: what does it mean for a promise of the form ``If $P$, then $Q$'' to be true when the condition $P$ does not happen?

Philo proposed an answer that startled the ancient world: an implication is true in every single case except when the premise $P$ is true and the outcome $Q$ is false. Under Philo's rule, whenever $P$ is false, the implication is automatically true---what modern logicians call \emph{material implication}, the rule that evaluates a conditional purely by the truth or falsehood of its parts, without requiring any real-world cause connecting them. Diodorus Cronus and the Stoic philosopher Chrysippus vehemently objected, arguing that an authentic conditional requires an intrinsic causal or logical bond. To their ears, calling ``If $2 + 2 = 5$, then the Moon is made of green cheese'' true seemed to mock human intuition. Bertrand Russell later dramatized the Megarian logic when challenged to prove that a false premise implies anything: if $1 = 0$, adding $1$ yields $2 = 1$; the collection containing Bertrand Russell and the Pope has two members, but $2 = 1$, so the collection has one member, proving that Russell and the Pope are one and the same person!

In 1879, Gottlob Frege and later Bertrand Russell and Alfred North Whitehead in \emph{Principia Mathematica} (1910) recognized that Philo's material implication was indispensable for formalizing mathematics. By decoupling mathematical proofs from subjective human feelings of ``relevance,'' mathematical logic could be executed as a rigorous, symbolic calculus. Without vacuous truth, modern set theory would instantly stall: we could not prove that the empty set is a subset of every set ($\emptyset \subseteq A$), nor could computer programs execute conditional branching instructions reliably.
\end{historybox}

In set theory, vacuous truth is the exact reason why the empty set $\emptyset$ is a subset of every set:
\[
\emptyset \subseteq A \quad \text{for every set } A.
\]
Recall the definition of subset: $S \subseteq A$ means that for every object $x$, if $x \in S$, then $x \in A$. When $S = \emptyset$, the hypothesis ``$x \in \emptyset$'' is false for every object in the mathematical universe. The condition is vacuously satisfied for every element, and the containment holds without exception.

\subsection{The trap of the converse and inverse}

Once you have an implication $P \implies Q$, there are three other conditional statements that can be formed by rearranging and negating its pieces:

\begin{enumerate}
\item \textbf{The Converse ($Q \implies P$)}: Swap the hypothesis and conclusion.
\item \textbf{The Inverse ($\neg P \implies \neg Q$)}: Negate both sides without swapping.
\item \textbf{The Contrapositive ($\neg Q \implies \neg P$)}: Swap both sides \emph{and} negate both sides.
\end{enumerate}

To see how these differ, let us test them on a simple arithmetic statement where $P$ is ``$x = 3$'' and $Q$ is ``$x^2 = 9$'', working with real numbers:
\begin{itemize}
\item \textbf{Original ($P \implies Q$)}: ``If $x = 3$, then $x^2 = 9$.'' \\
\emph{Status:} \textbf{True.} $3^2 = 9$.
\item \textbf{Converse ($Q \implies P$)}: ``If $x^2 = 9$, then $x = 3$.'' \\
\emph{Status:} \textbf{False!} If $x^2 = 9$, $x$ could equal $-3$.
\item \textbf{Inverse ($\neg P \implies \neg Q$)}: ``If $x \neq 3$, then $x^2 \neq 9$.'' \\
\emph{Status:} \textbf{False!} Setting $x = -3$ satisfies $x \neq 3$, yet $x^2 = 9$.
\item \textbf{Contrapositive ($\neg Q \implies \neg P$)}: ``If $x^2 \neq 9$, then $x \neq 3$.'' \\
\emph{Status:} \textbf{True.} If the square is not $9$, the number could not possibly have been $3$.
\end{itemize}

\begin{center}
\small
\begin{tabular}{lll}
\toprule
\textbf{Statement Form} & \textbf{Symbolic Logic} & \textbf{Truth Relation to $P \implies Q$} \\
\midrule
Original Implication & $P \implies Q$ & The baseline claim \\
Contrapositive & $\neg Q \implies \neg P$ & \textbf{Logically equivalent} (identical truth value) \\
Converse & $Q \implies P$ & \textbf{Independent statement} (may be true or false) \\
Inverse & $\neg P \implies \neg Q$ & \textbf{Equivalent to converse} (independent of original) \\
\bottomrule
\end{tabular}
\end{center}

Notice that the inverse ($\neg P \implies \neg Q$) is the contrapositive of the converse ($Q \implies P$). That is why the inverse and converse always share the exact same truth value. But neither of them has any logical obligation to agree with the original statement!

The square of conditionals in the diagram below summarizes these relationships: diagonal pairs share identical truth values, while horizontal and vertical moves lead to independent claims.

\begin{center}
\begin{tikzpicture}[
    font=\small,
    boxorig/.style={draw=sourceblue!80, fill=sourceblue!8, rounded corners=4pt, thick, text width=4.6cm, inner sep=6pt, align=center},
    boxconv/.style={draw=quotientorange!80, fill=quotientorange!8, rounded corners=4pt, thick, text width=4.6cm, inner sep=6pt, align=center},
    linkdash/.style={-{Latex[length=2.5mm]}, thick, trapred!80, dashed}
]

% Top-Left: Original
\node[boxorig] (orig) at (0, 3.4) {
    \textbf{\color{sourceblue!90!black}Original Implication}\\[2pt]
    $P \implies Q$\\[2pt]
    \scriptsize ``If $x = 3$, then $x^2 = 9$''\\\textbf{\color{targetgreen!80!black}Status: True}
};

% Top-Right: Converse
\node[boxconv] (conv) at (7.8, 3.4) {
    \textbf{\color{quotientorange!90!black}Converse (Swap sides)}\\[2pt]
    $Q \implies P$\\[2pt]
    \scriptsize ``If $x^2 = 9$, then $x = 3$''\\\textbf{\color{trapred!80!black}Status: False (counterexample $x=-3$)}
};

% Bottom-Left: Inverse
\node[boxconv] (inv) at (0, -0.8) {
    \textbf{\color{quotientorange!90!black}Inverse (Negate sides)}\\[2pt]
    $\neg P \implies \neg Q$\\[2pt]
    \scriptsize ``If $x$ is not $3$, then $x^2$ is not $9$''\\\textbf{\color{trapred!80!black}Status: False (counterexample $x=-3$)}
};

% Bottom-Right: Contrapositive
\node[boxorig] (contra) at (7.8, -0.8) {
    \textbf{\color{sourceblue!90!black}Contrapositive (Swap \& Negate)}\\[2pt]
    $\neg Q \implies \neg P$\\[2pt]
    \scriptsize ``If $x^2$ is not $9$, then $x$ is not $3$''\\\textbf{\color{targetgreen!80!black}Status: True}
};

% Horizontal dashed arrows (non-equivalence)
\draw[linkdash] (orig.east) -- node[above=2pt, font=\scriptsize\bfseries, text=trapred!90!black]{Swap sides} node[below=2pt, font=\tiny, text=trapred!80!black]{Not equivalent!} (conv.west);
\draw[linkdash] (inv.east) -- node[above=2pt, font=\scriptsize\bfseries, text=trapred!90!black]{Swap sides} node[below=2pt, font=\tiny, text=trapred!80!black]{Not equivalent!} (contra.west);

% Vertical dashed arrows (non-equivalence)
\draw[linkdash] (orig.south) -- node[left=3pt, font=\scriptsize\bfseries, text=trapred!90!black, align=right]{Negate\\sides} (inv.north);
\draw[linkdash] (conv.south) -- node[right=3pt, font=\scriptsize\bfseries, text=trapred!90!black, align=left]{Negate\\sides} (contra.north);

% Central equivalence banner
\node[draw=targetgreen!80, fill=targetgreen!12, rounded corners=4pt, thick, inner sep=5pt, align=center, font=\scriptsize] (center) at (3.9, 1.3) {
    \textbf{\color{targetgreen!90!black}Diagonal Equivalence}\\[2pt]
    $(P \implies Q) \iff (\neg Q \implies \neg P)$\\[2pt]
    $(Q \implies P) \iff (\neg P \implies \neg Q)$
};

\end{tikzpicture}
\end{center}

\begin{trapbox}
\textbf{The Converse Fallacy}: One of the most common errors in all of mathematics is confusing an implication with its converse. In casual conversation, people say ``If it is raining, the sidewalk is wet'' and loosely assume that a wet sidewalk implies rain. But a sprinkler could have caused it! Never assume that proving $P \implies Q$ gives you any license to use $Q \implies P$.
\end{trapbox}

What happens when an implication and its converse are \emph{both} true? In that case, the two statements are logically equivalent, and we write:
\[
P \iff Q \quad (\text{``}P \text{ if and only if } Q\text{''}, \text{ often abbreviated } P \text{ iff } Q).
\]
A statement of the form $P \iff Q$ is called a \textbf{biconditional}---a two-way promise. Whenever you see a theorem framed as ``$P$ if and only if $Q$,'' you are being given two separate mathematical tasks:
\begin{enumerate}
\item \textbf{Forward direction ($\implies$)}: Assume $P$, deduce $Q$.
\item \textbf{Backward direction ($\impliedby$)}: Assume $Q$, deduce $P$ (or prove its contrapositive: assume $\neg P$, deduce $\neg Q$).
\end{enumerate}

\begin{recoverybox}
\textbf{Pause and think:} State the converse of: ``If a number is divisible by $10$, then it is divisible by $5$.'' Is the converse true?

\emph{Answer:} The converse is: ``If a number is divisible by $5$, then it is divisible by $10$.'' This converse is false: for instance, $15$ is divisible by $5$, but not by $10$. The original implication holds, but the converse fails.
\end{recoverybox}

\subsection{Three standard proof strategies}

When proving an implication $P \implies Q$, there are three standard approaches:

\begin{center}
\begin{tikzpicture}[
    font=\small,
    plan/.style={draw=#1!70!black, fill=#1!6, rounded corners=3pt, thick, text width=12.2cm, inner sep=6pt}
]

% Plan 1: Direct Proof
\node[plan=targetgreen] (p1) at (0, 3.2) {
    \textbf{\color{targetgreen!90!black}1. Direct Proof} \quad \footnotesize\textbf{When to use:} When $P$ provides rich algebraic or structural ingredients.\\
    \scriptsize
    \textbf{Template:} ``Assume $P$ is true. Unpack definitions into algebraic formulas. Deduce successive truths until $Q$ is established.''
};

% Plan 2: Contrapositive
\node[plan=sourceblue] (p2) at (0, 1.0) {
    \textbf{\color{sourceblue!90!black}2. Contrapositive Proof ($\neg Q \implies \neg P$)} \quad \footnotesize\textbf{When to use:} When $\neg Q$ is more concrete than $P$.\\
    \scriptsize
    \textbf{Template:} ``Assume $Q$ is false ($\neg Q$). Unpack the negated definition. Deduce that $P$ must also be false ($\neg P$). Conclude $P \implies Q$.''
};

% Plan 3: Contradiction
\node[plan=trapred] (p3) at (0, -1.2) {
    \textbf{\color{trapred!90!black}3. Proof by Contradiction} \quad \footnotesize\textbf{When to use:} When assuming $\neg Q$ quickly leads to a contradiction with a known property.\\
    \scriptsize
    \textbf{Template:} ``Assume for contradiction that $P$ is true AND $Q$ is false ($P \land \neg Q$). Derive a logical contradiction ($R \land \neg R$). Conclude $Q$ must hold.''
};

\end{tikzpicture}
\end{center}

\subsection{Strategy 1: The direct route and the forward-backward method}

In a direct proof, you take premise $P$ as given and deduce conclusion $Q$ through a chain of lawful steps. But how do you find those steps when staring at a blank page?

Beginners often imagine that mathematicians discover proofs in a single forward pass, writing down line 1, then line 2, all the way to the end. In reality, finding a proof is often a two-directional process: the \textbf{forward-backward method}.

Before writing the final proof, split your scratchpad into two columns:
\begin{enumerate}
\item \textbf{The Forward Column (What do I know?)}: Start with hypothesis $P$. Unpack its definitions into concrete algebraic tools. Ask: \emph{``What facts can I extract immediately from this starting point?''}
\item \textbf{The Backward Column (What do I need?)}: Look at conclusion $Q$. Unpack what it would take to justify $Q$. Ask: \emph{``What statement, if I could prove it, would be sufficient to declare $Q$ true?''}
\item \textbf{Connecting the Two Sides}: Manipulate the forward equations until they match the form required by the backward target. Once the two ends connect, write down the final proof from start to finish, presenting a clean deductive argument.
\end{enumerate}

Let us see how this method works on a concrete example.

\begin{theorem}
Let $n \in \Z$. If $n$ is odd, then $5n + 3$ is even.
\end{theorem}

\begin{proof}[Scratchwork and Discovery]
\textbf{Forward:} The hypothesis is that $n$ is odd. Unpacking the definition of odd from Chapter~6 gives:
\[
n = 2k + 1 \quad \text{for some integer } k.
\]
\textbf{Backward:} The target conclusion is that $5n + 3$ is even. Unpacking the definition of even tells us our destination: we must show that
\[
5n + 3 = 2m \quad \text{for some integer } m.
\]
\textbf{Connecting the bridge:} We substitute the forward expression $n = 2k + 1$ into $5n + 3$:
\[
5n + 3 = 5(2k + 1) + 3 = 10k + 5 + 3 = 10k + 8 = 2(5k + 4).
\]
Because $k$ is an integer, $5k + 4$ is an integer. Setting $m = 5k + 4$ delivers the exact form $2m$ demanded by the definition of an even integer.
\end{proof}

The diagram below visualizes this strategy: forward unpacking and backward target-setting meet once the expression is factored into the required form.

\begin{center}
\begin{tikzpicture}[
    font=\small,
    front/.style={draw=#1!80, fill=#1!8, rounded corners=5pt, thick, text width=3.3cm, inner sep=6pt, align=center},
    steparrow/.style={-{Latex[length=2.5mm]}, very thick}
]

% Left: Forward Column
\node[front=sourceblue] (fwd) at (1.1, 0) {
    \textbf{\color{sourceblue!90!black}Forward Steps}\\[2pt]
    \scriptsize
    \textbf{1. Hypothesis:}\\
    $n$ is an odd integer\\[3pt]
    $\Downarrow$ \emph{unpack definition}\\[3pt]
    $n = 2k + 1$ \ ($k \in \Z$)\\[3pt]
    $\Downarrow$ \emph{substitute into target}\\[3pt]
    $5n + 3 = 5(2k + 1) + 3$\\
    $= 10k + 8$
};

% Center: The Bridge
\node[front=structurepurple, text width=4.3cm] (br) at (6.6, 0) {
    \textbf{\color{structurepurple!90!black}Connecting the Sides}\\[3pt]
    \scriptsize
    Factor target into required form:\\[2pt]
    $10k + 8 = 2(5k + 4)$\\[4pt]
    Identify witness integer:\\[2pt]
    $m = 5k + 4$\\[4pt]
    Integer closure check:\\[2pt]
    $5k + 4 \in \Z$ \quad \textbf{\color{targetgreen!80!black}\checkmark}\\[3pt]
    Exact match to backward need!
};

% Right: Backward Column
\node[front=targetgreen] (bwd) at (12.1, 0) {
    \textbf{\color{targetgreen!90!black}Backward Target}\\[2pt]
    \scriptsize
    \textbf{Target Form:}\\
    $5n + 3 = 2m$ \ ($m \in \Z$)\\[3pt]
    $\Uparrow$ \emph{package definition}\\[3pt]
    Must show $5n + 3$ is a multiple of $2$\\[3pt]
    $\Uparrow$ \emph{target goal}\\[3pt]
    \textbf{Conclusion:}\\
    $5n + 3$ is even
};

% Connecting Horizontal Arrows
\draw[steparrow, sourceblue!80] (fwd.east) -- (br.west);
\draw[steparrow, targetgreen!80] (bwd.west) -- (br.east);

\end{tikzpicture}
\end{center}

Now we write the polished formal proof:

\begin{proof}
Assume $n$ is an odd integer.

By definition of an odd integer, there exists an integer $k$ such that $n = 2k + 1$. Substituting this expression into $5n + 3$, we have:
\[
5n + 3 = 5(2k + 1) + 3 = 10k + 5 + 3 = 10k + 8 = 2(5k + 4).
\]
Let $m = 5k + 4$. Since $k \in \Z$, the quantity $m$ is an integer by closure of the integers under addition and multiplication. 

Thus, $5n + 3 = 2m$ for an integer $m$. By definition, $5n + 3$ is even.
\end{proof}

Every step is accounted for: we unpacked the definition of odd, carried out the algebra, checked integer closure, and verified the definition of an even integer.

\subsection{Strategy 2: Contrapositive proof}

Direct proof works well when the hypothesis $P$ gives concrete algebraic formulas to manipulate. But what happens when the hypothesis is awkward to unpack?

Consider the following innocent-looking theorem:
\begin{theorem}
Let $n \in \Z$. If $n^2$ is even, then $n$ is even.
\end{theorem}

Let us attempt a direct proof. We assume $n^2$ is even. Unpacking the definition gives:
\[
n^2 = 2k \quad \text{for some integer } k.
\]
Now we want to deduce that $n$ is even. To isolate $n$, we take square roots:
\[
n = \sqrt{2k}.
\]
Here the direct approach stalls. Is $\sqrt{2k}$ an integer? Squaring an integer is straightforward, but taking square roots leaves divisibility unclear.

This is the exact moment to pivot to the \textbf{contrapositive}:
\[
\neg Q \implies \neg P.
\]
Recall that for any statements $P$ and $Q$, the implication $P \implies Q$ is logically equivalent to its contrapositive:
\[
(P \implies Q) \iff (\neg Q \implies \neg P).
\]
Look back at the Euler diagram. To say that $P$ is inside $Q$ ($P \subseteq Q$) is geometrically identical to saying that if you stand outside $Q$, you are guaranteed to be standing outside $P$ ($Q^c \subseteq P^c$). If you are outside the big circle, you cannot possibly be inside the small circle! If a bedroom is inside an apartment, someone standing outside the apartment is guaranteed not to be in the bedroom.

Now let us apply the contrapositive to our theorem. The conclusion is ``$n$ is even,'' so its negation $\neg Q$ is ``$n$ is not even'' (meaning $n$ is odd). The hypothesis is ``$n^2$ is even,'' so its negation $\neg P$ is ``$n^2$ is odd.'' The contrapositive claim is:
\begin{center}
\itshape ``If $n$ is odd, then $n^2$ is odd.''
\end{center}
Instead of starting with $n^2$ and taking square roots, we start with $n$ and square it. The algebraic direction is straightforward.

\begin{proof}
We prove the theorem by contrapositive: if $n$ is not even, then $n^2$ is not even.

For integers, an integer that is not even is odd. Assume $n$ is odd. By definition, there exists an integer $k$ such that $n = 2k + 1$. Squaring both sides:
\[
n^2 = (2k + 1)^2 = 4k^2 + 4k + 1 = 2(2k^2 + 2k) + 1.
\]
Because $k$ is an integer, $m = 2k^2 + 2k$ is an integer by closure under addition and multiplication. Thus $n^2 = 2m + 1$, which means $n^2$ is odd.

We have established that if $n$ is odd, then $n^2$ is odd. Because an implication is logically equivalent to its contrapositive, it follows that if $n^2$ is even, then $n$ must be even.
\end{proof}

\subsection{Choosing a proof strategy: direct, contrapositive, or contradiction}

When proving an implication, which strategy should you choose? Here are standard guidelines:

\begin{itemize}
\item \textbf{Use Direct Proof} when hypothesis $P$ provides clear algebraic formulas (such as $n$ is odd, $a \mid b$, or $x, y \in \Q$) that you can unpack and substitute directly into the conclusion. If the forward path moves naturally toward $Q$, direct proof is the cleanest and most transparent route.
\item \textbf{Use Contrapositive} when:
\begin{itemize}
\item The conclusion $Q$ is simpler to unpack than hypothesis $P$ (for example, starting from $n$ is far easier than starting from $n^2$).
\item The conclusion contains a negative assertion, such as ``$x$ is irrational'' or ``$x \neq 0$.'' Negating $\neg Q$ transforms a negative condition into a positive algebraic equation ($x \in \Q$ or $x = 0$) that gives you immediate formulas to manipulate.
\end{itemize}
\item \textbf{Use Proof by Contradiction} when neither direct nor contrapositive paths offer a foothold. Direct and contrapositive proofs start from what is given and deduce consequences forward. Proof by contradiction assumes both the hypothesis and the negation of the conclusion ($P \land \neg Q$), reasoning from there until reaching a contradiction ($R \land \neg R$).
\end{itemize}

Contradiction is especially useful when a theorem makes a negative claim with no algebraic hypothesis to unpack---such as proving that $\sqrt{2}$ is irrational, or that the primes never end.

In Chapter~8 (\emph{Proof by Contradiction: Power and Danger}), we examine this method in detail: how assuming the opposite uncovers contradictions, classic examples from ancient Greece to modern analysis, and how to avoid common pitfalls.

\subsection{Exercises}

\subsubsection*{Part A: Warm-Up Mental Checks}
\begin{enumerate}
\item Evaluate the truth value of each implication:
\begin{enumerate}
\item If $2 + 2 = 5$, then $3 \times 3 = 9$.
\item If $2 + 2 = 4$, then $3 \times 3 = 10$.
\item If $2 + 2 = 5$, then $3 \times 3 = 10$.
\end{enumerate}
\item For each statement, write down its converse, inverse, and contrapositive:
\begin{enumerate}
\item If an integer $n$ is a multiple of $8$, then $n$ is a multiple of $4$.
\item If a real number $x$ satisfies $x > 3$, then $x^2 > 9$.
\end{enumerate}
\item Identify the logical error in the following argument: ``Every differentiable function is continuous. Since $f(x) = |x|$ is continuous at $x = 0$, it must be differentiable at $x = 0$.''
\end{enumerate}

\subsubsection*{Part B: Routine Practice}
\begin{enumerate}[resume]
\item State the contrapositive and the negation of each implication:
\begin{enumerate}
\item If $x > 2$, then $x^2 > 4$.
\item If a triangle is equilateral, then it is equiangular.
\item If an integer $n$ is divisible by $6$, then $n$ is divisible by $2$ and by $3$.
\end{enumerate}
\item Prove each statement using a direct proof. Explicitly label where you unpack definitions and where you package your conclusion:
\begin{enumerate}
\item For $n \in \Z$, if $n$ is even, then $3n^2 - 5n + 6$ is even.
\item For $a, b, c \in \Z$, if $a \mid b$ and $a \mid c$, then $a \mid (2b + 3c)$.
\end{enumerate}
\end{enumerate}

\subsubsection*{Part C: Contrapositive Proofs}
\begin{enumerate}[resume]
\item Prove each statement by contrapositive:
\begin{enumerate}
\item For $n \in \Z$, if $3n + 2$ is odd, then $n$ is odd.
\item For real numbers $x, y$, if $x + y > 10$, then $x > 5$ or $y > 5$.
\item For an integer $n$, if $n^2 - 1$ is not divisible by $4$, then $n$ is even.
\end{enumerate}
\end{enumerate}

\subsubsection*{Part D: Deeper Challenges}
\begin{enumerate}[resume]
\item \textbf{A Biconditional Theorem:} Let $n \in \Z$. Prove that $n$ is even if and only if $n^2$ is even. (Explicitly present both the forward direction $\implies$ and the backward direction $\impliedby$, noting where the direct route and the contrapositive route each apply.)
\item \textbf{Rational and Irrational Numbers:} Recall that a real number $r$ is rational ($r \in \Q$) if $r = a/b$ for integers $a, b$ with $b \neq 0$, and irrational otherwise. Prove that if $x$ is irrational and $r$ is a non-zero rational number, then $r \cdot x$ is irrational. (Why does the contrapositive make working with irrational numbers more direct?)
\item If $x$ is irrational, prove that $-x$ is irrational.
\end{enumerate}

\newpage
